\documentclass[a4paper,fontsize=14bp  % c пакетом scrextend
%twoside,
%draft
]{extreport}

\input{vkr_identification}  % персональные данные
\usepackage{tex/PnuPhysStyle}      % не трогать

% Любые добавления новых определений и подключаемых пакетов включается здесь

%%%%%%%%%%%%%%%%%%%%%%%%%%%%%%%%%%%%%%%%%%%%%%%%%%%%%%%%%%%%%%%%%%%%%%%%%%%%%%% 
\begin{document}

\input{tex/vkr_title_page}  % не трогать
\input{tex/vkr_dopusk_page} %% не трогать
\input{tex/vkr_problem_page} % не трогать (при крайней необходимости можно)
\input{vkr_referat}      % Реферат - заполняем

\newpage
\renewcommand{\contentsname}{\MakeUppercase{Содержание}}
\tableofcontents %  автоматическое оглавление

	\input{vkr_introduction}  % Введение  - заполняем
	\input{vkr_chapter_1}  % Глава 1.  - заполняем
	\input{vkr_chapter_2}  % Глава 2.  - заполняем
	\input{vkr_chapter_3}  % Глава 3.  - заполняем
	\input{vkr_conclusion} % Заключение. - заполняем
%	\bibliography{vkr_bibliography} % Список литературы.
\linespread{1.1}
	\printbibliography[title={Список использованных источников},heading=bibintoc]

    \addtocounter{appxpage}{-\value{page}}
	\input{vkr_appendix}   % Приложение. - заполняем
    \addtocounter{appxpage}{\value{page}} 
    \refstepcounter{appxpage}
    \label{lastaapx}

	\input{tex/vkr_toc_footer}   % доп текст в содержании, не трогать (при крайней необходимости можно)

\end{document}